\documentclass[10pt,a4paper]{article}
\usepackage[margin=18mm]{geometry}
\usepackage[T1]{fontenc}
\usepackage{lmodern}
\usepackage{microtype}
\usepackage[hidelinks]{hyperref}
\usepackage{parskip}
\setlength{\parindent}{0pt}
\setlength{\parskip}{1pt plus 0.5pt minus 0.5pt}
\pagestyle{plain}

\newcommand{\sourceheading}[2]{%
  \section*{#1}
  \textbf{#2}\par\medskip
}

\begin{document}

\begin{center}
  {\Large\bfseries LN905 Reading into Writing Practice}\par
  \vspace{4pt}
  {\large Climate Change - End-to-End Paper B Unit}
\end{center}

\section*{Question}

\textbf{To what extent can adaptation prevent climate change from widening global inequality?}

Write an essay of approximately 600 words. Use all three extracts to develop your own qualified answer. The extracts are exam-style adaptations of the cited academic papers. Their wording has been condensed and paraphrased while preserving the studies' central claims, evidence and limitations.

\newpage

\sourceheading{Extract 1}{Adapted from Diffenbaugh and Burke (2019), ``Global Warming Has Increased Global Economic Inequality''}

Economic inequality between countries has declined over recent decades, but this overall trend does not reveal whether climate change altered the speed or distribution of that convergence. Diffenbaugh and Burke investigate whether anthropogenic warming changed national economic growth between 1961 and 2010. Their question is historical and counterfactual: how different would countries' output have been if human activity had not raised temperatures?

The researchers combine two kinds of modelling evidence. First, climate-model simulations estimate each country's temperature path with and without anthropogenic forcing. Second, an empirically estimated relationship links annual temperature fluctuations to economic growth. That relationship is nonlinear: additional warmth can improve growth in relatively cool countries but tends to reduce growth in countries that are already warm. By repeatedly combining uncertainty in both components, the authors construct more than 20,000 possible no-warming histories and compare these with observed economic outcomes.

The estimated distribution is strongly unequal. Across the four poorest population-weighted deciles, per capita gross domestic product in 2010 is estimated to have been 17-31\% lower than it would have been without global warming. The ratio between the richest and poorest deciles is estimated to be 25\% larger, and the authors place roughly 90\% probability on warming having slowed the decline in between-country inequality. Countries in tropical and subtropical regions usually experienced losses, while some cooler and richer countries may have benefited. Among countries with very low cumulative emissions, all eighteen in one comparison had negative median impacts; among the highest historical emitters, fourteen of nineteen had positive median impacts.

The mechanism is not simply that poor countries have fewer resources. Many poor countries are also located in already warm climates, where further temperature increases move conditions farther from the estimated economic optimum. Because rich countries are more often temperate or cool, the same physical warming can produce different growth effects. The authors therefore argue that warming generated largely by high-emitting economies has reinforced an inequality rooted partly in earlier differences in fossil-fuel use.

This geographical overlap means that countries least responsible for cumulative emissions can be among those most exposed to climate-related growth losses.

The comparison also separates the effect of warming from the many other forces shaping development. Several poor economies still grew faster than richer economies during the study period, so the finding is not that climate change stopped convergence altogether. Rather, the counterfactual estimates suggest that convergence would probably have proceeded faster in a world without anthropogenic warming. This distinction matters for evaluating policy: an adaptation programme could improve a country's welfare relative to continued inaction while still leaving it behind the position it might have reached without past warming. Preventing further divergence is therefore a narrower objective than compensating for accumulated inequality.

The authors test alternative climate-model simulations, economic specifications and assumptions about each country's no-warming temperature. Although the size of estimated national impacts changes across these choices, the unequal global pattern remains. The evidence is thus more secure for the direction and distribution of the historical effect than for a precise monetary loss in any single country.

The study supports a strong claim about the likely direction of historical between-country inequality, but its exact magnitudes remain model-dependent. The economic response function is estimated from historical temperature variation and aggregate national output; it cannot identify every pathway through agriculture, health or labour productivity, and it does not measure inequality within countries. Some temperate countries also have uncertain effects. The analysis shows that warming probably widened international economic gaps; it does not establish that adaptation policies alone can reverse those accumulated losses.

{\itshape Source: Proceedings of the National Academy of Sciences, 116(20), 9808-9813.\par DOI: \href{https://doi.org/10.1073/pnas.1816020116}{10.1073/pnas.1816020116}.}

\newpage

\sourceheading{Extract 2}{Adapted from Carleton et al. (2022), ``Valuing the Global Mortality Consequences of Climate Change Accounting for Adaptation Costs and Benefits''}

Carleton and colleagues examine one important channel through which climate change may create unequal losses: temperature-related mortality. They aim to combine the global scale of long-run climate models with the empirical detail of local mortality data. Their dataset contains annual subnational mortality statistics from forty countries covering 38\% of the world's population. The analysis estimates age-specific relationships between daily temperature and death rates, then allows those relationships to vary with long-run climate and income. The resulting model is used to project outcomes for 24,378 regions under alternative climate and socioeconomic scenarios.

Extreme cold and extreme heat are associated with higher mortality, especially among people over sixty-four. However, the response varies substantially across locations. Higher income and experience of a warmer climate are associated with lower sensitivity to extreme temperatures, which the authors interpret as evidence of income-related protection and adaptation. Under a high-emissions scenario, their mean projection for 2100 is 73 additional deaths per 100,000 people. Ignoring both future income growth and adaptation would raise the estimate to 221 deaths per 100,000, approximately three times as large.

Adaptation therefore has substantial potential, but it does not distribute risk equally or eliminate it. In Accra, Ghana, the model projects about a 17\% increase in the current mortality rate by 2100 as very hot days become more frequent. In Berlin, Germany, it projects a 15\% decline because fewer cold days offset heat-related deaths. More generally, poorer locations are projected to experience more of the climate burden as deaths, whereas richer locations avoid more deaths but pay for protective investments. Adaptation capacity can reduce physical harm, yet the resources required to adapt become part of the total social cost.

The global average therefore conceals two sources of inequality. Exposure differs because warming shifts the number of dangerously hot and cold days in different climates. Protection also differs because locations enter the future with unequal income, infrastructure and experience of extreme heat. Economic growth is projected to strengthen protection in many places, but the poorest and hottest regions can remain highly exposed even after the model allows adaptation. A reduction in aggregate deaths does not by itself show that the gap between regions has narrowed.

This distinction is important when assessing whether adaptation is sufficient. Measures such as cooling, information, health services and changes in work or behaviour can reduce mortality, but their availability and cost depend on local resources. If high-income regions can purchase more protection while low-income regions retain larger residual losses, adaptation may lower total harm without preventing climate change from widening inequality. The distribution of both avoided deaths and adaptation expenditure must therefore be considered.

To include those costs, the authors infer adaptation expenditure from observed differences in temperature sensitivity. Their revealed-preference model assumes that representative agents invest until the marginal value of avoided deaths equals the marginal cost of further adaptation. Including projected deaths and inferred adaptation costs, the mean full mortality risk under high emissions is valued at about 3.2\% of global GDP in 2100. The interquartile range, from -5.4\% to 9.1\%, shows that substantial uncertainty surrounds the estimate.

Several limits matter for using this evidence. Future mortality depends on emissions, income, demographics and technologies that cannot be known with certainty. The observed relationship between income or long-run climate and mortality sensitivity is cross-sectional, so the authors treat it as associational rather than definitively causal. Their cost model also assumes optimizing agents and relatively frictionless markets. Finally, mortality is only one consequence of climate change; storms, conflict, migration and many other outcomes are outside this partial calculation. The findings nevertheless indicate that adaptation can greatly reduce harm while leaving large, uneven residual risks.

{\itshape Source: Quarterly Journal of Economics, 137(4), 2037-2105.\par DOI: \href{https://doi.org/10.1093/qje/qjac020}{10.1093/qje/qjac020}.}

\newpage

\sourceheading{Extract 3}{Adapted from Berrang-Ford et al. (2021), ``A Systematic Global Stocktake of Evidence on Human Adaptation to Climate Change''}

Research frequently models what adaptation could achieve, but policy also depends on what governments, firms, communities and households have actually implemented. Berrang-Ford and a large international team review empirical studies of observed human responses to climate change. Machine-learning methods were used to screen 48,816 peer-reviewed articles published between 2013 and 2019. A network of 126 researchers then coded the 1,682 articles that met the inclusion criteria, asking who was responding, which hazards motivated action, what form the response took, whether it reduced climate risk and whether it was transformational.

The literature documents adaptation in every major world region, with the largest shares of articles concerning Asia and Africa. Nearly two-thirds focus on food and agriculture, reflecting both that sector's climate sensitivity and the visibility of household and producer responses. Reported actions include changing agricultural practices, improving infrastructure, using technological or behavioural measures and adjusting institutions. This broad evidence base confirms that adaptation is not merely theoretical: people and organizations are already responding to changing climatic conditions.

However, the pattern of action is less reassuring than the quantity of published research might suggest. The documented responses are predominantly local, fragmented and incremental. They usually extend existing practices rather than changing the values, institutions or structures that create vulnerability. Evidence of coordinated or transformational adaptation is limited, even in settings facing high present or projected risks. More importantly, the literature contains negligible evidence that the reported actions have actually reduced climate risk. An implemented response and an effective response are therefore not equivalent.

This distinction changes how adaptation should be judged. A farmer may adopt a new crop, a city may publish a heat plan, or a household may purchase cooling equipment, yet these actions do not automatically protect the most vulnerable groups or remain effective as warming intensifies. Small, local adjustments can be useful, but they may be overwhelmed by larger hazards, constrained by finance and institutions, or reproduce existing inequalities. The authors identify priorities including evaluating outcomes, understanding hard and soft limits, enabling individuals and civil society to act, studying private-sector responses and including places and scholarship that are missing from the current evidence base.

The stocktake also reveals a measurement problem. Studies commonly describe that an action occurred, but far fewer evaluate whether exposure, sensitivity or vulnerability changed afterwards. Evidence about distribution is thinner still: an average improvement can coexist with continuing disadvantages for low-income households, women, marginalized communities or regions with limited political capacity. Without outcome and equity measures, it is difficult to tell whether adaptation has prevented inequality, shifted risk to another group or merely delayed damage.

Scale and coordination are therefore central. Local experiments may offer protection, yet hazards cross borders while finance and decision-making power remain uneven. Responses that remain isolated can leave structural causes of vulnerability untouched. The review does not imply that every effective intervention must be transformational, but it suggests that incremental actions need to be connected, financed and evaluated if they are to address inequality rather than only isolated symptoms.

The review also has important boundaries. It maps what peer-reviewed studies document, not every adaptation occurring in the world. Activities that are unpublished, described in non-English or non-academic sources, or difficult to observe may be underrepresented. Publication counts do not measure investment or effectiveness, and the scarcity of risk-reduction evidence does not prove that adaptation has failed; it may instead reflect weak monitoring and evaluation. Even so, the stocktake challenges the assumption that recognising a risk will naturally produce adaptation at sufficient scale. Preventing climate change from deepening inequality requires not only more responses but evidence that they reduce vulnerability, reach disadvantaged groups and remain effective under greater warming.

{\itshape Source: Nature Climate Change, 11, 989-1000.\par DOI: \href{https://doi.org/10.1038/s41558-021-01170-y}{10.1038/s41558-021-01170-y}.}

\end{document}
